%% 第13回編集委員会（物理第1巻）
%% 物理シリーズ　第1巻　力学　巻末（こたえ・記号カード・物性値・定数・単位・なぜ？の道すじ・チェック表・つぎの巻へ・さくいん・奥付）
%% 章を書いたあとで中身を足す。いまは形と仮の中身（物性値の表の数は仮。出典を確かめてから確定する）。
%% 見出し（\bkbacktitle）の名前は、もくじ（ph1-toc.tex）の巻末の項目の名前と同じにする（tools/check_pages.py p1 が照らし合わせる）。
\makeatletter
\def\phbkhead#1{\par\addvspace{4mm}\Needspace{6\baselineskip}\noindent{\sffamily\bfseries\large\color{bkDeep}#1}\par\nobreak\vspace{2mm}}
\newcommand\ph@tblstyle{row{1}={bg=bkVol!14, font=\sffamily\bfseries\footnotesize}, hlines={0.4pt,bkLine}, rowsep=1.4pt, colsep=3.5pt}
%% ---------------- こたえ ----------------
\bkAnswerPages
\par\vspace{6mm}
\begin{tcolorbox}[bk base, colframe=bkTheme!60, colback=bkTheme!5, boxrule=0.8pt, arc=2mm, left=4mm, right=4mm, top=2.5mm, bottom=3mm]
{\sffamily\bfseries\color{bkDeep}答え合わせのしかた}\par\vspace{1.5mm}
\begin{enumerate}[label={\bkcircled{\arabic*}}, leftmargin=2.4\zw, itemsep=1pt, topsep=0pt]
\item まず自分の答えを全部書いてから、ここを見よう。
\item 数だけでなく、単位と、桁（10倍・100倍のちがい）を確かめよう。物理の答えは、単位がちがえば別の量です。
\item 有効数字がちがうだけで、最後の桁が1つずれることがあります。それはまちがいではありません。
\item ちがっていたら、同じ形の例題にもどって、どの段（何を無視するか・式を立てる・解く）でちがったかをさがそう。
\end{enumerate}
\end{tcolorbox}

%% ---------------- 記号カード ----------------
%% 数学シリーズの記号で、この巻でよく使うもの（この巻の章の \bkkigocard は自動でうしろに集まる）
\bkKigoPages{%
\bkkigorow{$\Delta$}{デルタ}{変化量。「あと$-$まえ」}{2-3}%
\bkkigorow{$\theta$}{シータ}{角の大きさ（微分・積分ではラジアン）}{4-2}%
\bkkigorow{$\omega$}{オメガ}{角速度。1秒に回る角（rad/s）}{4-8}%
\bkkigorow{$\vec{a}\cdot\vec{b}$}{内積}{どれだけ同じ向きか。仕事 $W=\vec{F}\cdot\vec{s}$}{4-12}%
\bkkigorow{$\boldsymbol{a}$}{太字のベクトル}{専門書のベクトルの書き方。$\vec{a}$ と同じ}{4-13}%
\bkkigorow{$\vec{a}\times\vec{b}$}{外積}{回す働き。トルク $\vec{r}\times\vec{F}$}{4-14}%
\bkkigorow{$\dfrac{\dd x}{\dd t}$、\ $\dot{x}$}{ディーエックス・ディーティー、エックスドット}{時間で微分。速度}{5-5}%
\bkkigorow{$\displaystyle\int$}{インテグラル}{細かく分けてたす（積分）}{5-10}%
\bkkigorow{$\dfrac{\partial f}{\partial x}$}{ラウンドディー}{偏微分。ほかの変数を止めて $x$ だけ動かす}{8-2}%
\bkkigorow{$\mathrm{grad}$、\ $\nabla$}{グラディエント、ナブラ}{いちばん急な向き（勾配）}{8-4}%
\bkkigorow{$\mathrm{div}$}{ダイバージェンス}{わき出し（発散）}{8-9}%
\bkkigorow{$\mathrm{rot}$}{ローテーション}{うず（回転）}{8-10}%
}

\clearpage
\phbkhead{分野で記号がぶつかるもの}
{\small 物理と工学では、使える文字が足りないので、同じ文字がちがう量を表すことがあります。この本では次のように使い分けます。どの量かは、章と単位で見分けよう。（第13回の仮の表。章を書きながら足す）}\par\vspace{2mm}
{\footnotesize
\begin{longtblr}{colspec={Q[c,m,wd=11mm] X[l,m] X[l,m] Q[l,m,wd=33mm]}, width=\linewidth, rowhead=1, \ph@tblstyle}
記号 & この巻での意味 & ほかの章・巻での意味 & 見分け方 \\
$I$ & 慣性モーメント（力-6）、断面二次モーメント（力-10） & 電流（電-4） & 単位：kg·m$^2$、m$^4$、A \\
$p$ & 圧力（力-12）、運動量（力-3） & ― & 単位：Pa、kg·m/s。運動量のベクトルは $\boldsymbol{p}$ \\
$\sigma$ & 応力（力-8） & 標準偏差（3-14）、導電率（電-4） & 単位：Pa \\
$\varepsilon$ & ひずみ（力-8） & 誘電率（電-3） & ひずみは単位なし \\
$\tau$ & せん断応力（力-8） & 時定数（6-4） & 単位：Pa、s \\
$\mu$ & 摩擦係数（力-2）、粘度（力-14） & 透磁率（電-7） & 摩擦係数は単位なし、粘度は Pa·s \\
$\nu$ & ポアソン比（力-8）、動粘度（力-14） & 振動数（第3巻） & ポアソン比は単位なし、動粘度は m$^2$/s \\
$E$ & ヤング率（力-8）、エネルギー（力-4） & 電場（電-1） & 単位：Pa、J、V/m \\
$G$ & 横弾性係数（力-8）、万有引力定数 & ― & 単位：Pa、N·m$^2$/kg$^2$ \\
$k$ & ばね定数（力-2） & 波数（波-5）、ボルツマン定数（熱-14） & 単位：N/m、rad/m、J/K \\
$Q$ & 体積流量（力-13） & 熱量（熱-1）、電荷（電-1） & 単位：m$^3$/s、J、C \\
$P$ & 仕事率・動力（力-4）、荷重（力-11） & 電力（電-11） & 単位：W、N \\
$\rho$ & 密度（力-12） & 抵抗率（電-4） & 単位：kg/m$^3$、Ω·m \\
$T$ & 周期（力-1）、張力（力-2）、トルク & 温度（熱-1） & 単位：s、N、N·m、K \\
$h$ & 高さ（力-4） & 比エンタルピー（熱-4）、プランク定数（原-2） & 単位：m、J/kg、J·s \\
\end{longtblr}\par}

%% ---------------- 物性値・定数・単位 ----------------
\bkbacktitle{物性値・定数・単位}
\begin{tcolorbox}[bk base, colframe=bkRed!70!black, colback=bkRedBg, boxrule=0.8pt, arc=2mm, left=3mm, right=3mm, top=1.5mm, bottom=1.5mm, unbreakable]
\small 第13回の仮の表です。値と出典は、章を書くときに理科年表・機械工学便覧・JIS などで確かめて直します（「要確認」の印）。
\end{tcolorbox}
\phbkhead{1　材料の定数（室温）}
{\footnotesize
\begin{longtblr}{colspec={Q[l,m,wd=27mm] Q[c,m,wd=17mm] Q[c,m,wd=13mm] Q[c,m,wd=16mm] Q[c,m,wd=21mm] X[l,m]}, width=\linewidth, rowhead=1, \ph@tblstyle}
材料 & ヤング率 $E$ & ポアソン比 $\nu$ & 密度 $\rho$（kg/m$^3$） & 降伏応力など（MPa） & 出典 \\
鋼（一般構造用 SS400） & 約\qty{205}{GPa} & 約0.3 & 約7850 & 245以上（降伏点） & JIS G 3101 ほか（要確認） \\
アルミニウム合金（A6061-T6） & 約\qty{69}{GPa} & 約0.33 & 約2700 & 約270（耐力） & 要確認 \\
銅（タフピッチ銅、焼なまし） & 約\qty{117}{GPa} & 約0.34 & 約8900 & 約70（耐力） & 要確認 \\
コンクリート（普通） & 約\qty{23}{GPa} & 約0.2 & 約2300 & 圧縮強さ約24（引張はその約10分の1） & 要確認 \\
木材（すぎ、繊維の向き） & 約\qty{7}{GPa} & 約0.4 & 約380 & 圧縮強さ約35（降伏点はない） & 要確認 \\
\end{longtblr}\par}
{\footnotesize\color{bkGray}耐力：はっきりした降伏点がない金属で、0.2\,\% の永久ひずみが残る応力（力-8）。\par}

\phbkhead{2　流体の物性}
{\footnotesize
\begin{longtblr}{colspec={Q[l,m,wd=31mm] Q[c,m,wd=19mm] Q[c,m,wd=23mm] Q[c,m,wd=25mm] X[l,m]}, width=\linewidth, rowhead=1, \ph@tblstyle}
流体（温度・圧力） & 密度 $\rho$（kg/m$^3$） & 粘度 $\mu$（mPa·s） & 動粘度 $\nu$（mm$^2$/s） & 出典 \\
水（\qty{20}{\celsius}、1気圧） & 998.2 & 1.002 & 1.004 & 理科年表（要確認） \\
海水（\qty{15}{\celsius}） & 約1025 & 約1.2 & 約1.2 & 要確認 \\
空気（\qty{20}{\celsius}、1気圧） & 1.204 & 0.0181 & 15.0 & 理科年表（要確認） \\
作動油（ISO VG32、\qty{40}{\celsius}） & 約860 & 約28 & 32（決め方による） & ISO 3448（要確認） \\
\end{longtblr}\par}
{\footnotesize\color{bkGray}動粘度＝粘度÷密度（力-14）。$\qty{1}{mm^2/s}=\qty{1e-6}{m^2/s}$、$\qty{1}{mPa.s}=\qty{1e-3}{Pa.s}$。\par}

\phbkhead{3　この巻で使う定数}
{\footnotesize
\begin{longtblr}{colspec={Q[l,m,wd=40mm] Q[l,m,wd=52mm] X[l,m]}, width=\linewidth, rowhead=1, \ph@tblstyle}
名前 & 値 & 出典・注意 \\
標準の重力加速度 $g$ & \qty{9.80665}{m/s^2}（決めた値） & 本文では約\qty{9.8}{m/s^2}（第13回で決める） \\
標準大気圧 & \qty{101325}{Pa}（決めた値） & 約\qty{101}{kPa}、\qty{1013}{hPa} \\
万有引力定数 $G$ & \qty{6.67430e-11}{N.m^2/kg^2} & CODATA 2018 \\
地球の質量・半径 & \qty{5.97e24}{kg}、約\qty{6371}{km}（平均） & 要確認 \\
水の密度（\qty{4}{\celsius}） & 約\qty{1000}{kg/m^3}（999.97） & 理科年表（要確認） \\
\end{longtblr}\par}

\phbkhead{4　単位と換算}
{\footnotesize
\begin{longtblr}{colspec={Q[l,m,wd=30mm] Q[l,m,wd=48mm] X[l,m]}, width=\linewidth, rowhead=1, \ph@tblstyle}
量 & SI の単位 & ほかの単位との関係（単位つきの1で換算する。2-10） \\
力 & N（ニュートン）＝kg·m/s$^2$ & 1\,kgf＝\qty{9.80665}{N} \\
圧力・応力 & Pa（パスカル）＝N/m$^2$ & 1\,MPa＝1\,N/mm$^2$、1\,bar＝\qty{100}{kPa}、1\,kgf/cm$^2$＝\qty{98.0665}{kPa}、1\,mmHg≒\qty{133.3}{Pa}、1気圧＝\qty{101.325}{kPa} \\
エネルギー・仕事 & J（ジュール）＝N·m & 1\,kWh＝\qty{3.6}{MJ}、1\,cal＝\qty{4.184}{J} \\
仕事率・動力 & W（ワット）＝J/s & 1\,PS（仏馬力）≒\qty{735.5}{W} \\
トルク & N·m & 仕事の J と同じ次元だが、トルクは J と書かない \\
角速度・回転数 & rad/s & \qty{1}{rpm}（1分間の回転数）＝$\dfrac{2\pi}{60}$\,rad/s \\
粘度・動粘度 & Pa·s、m$^2$/s & 1\,P（ポアズ）＝\qty{0.1}{Pa.s}、1\,St（ストークス）＝\qty{1e-4}{m^2/s} \\
\end{longtblr}\par}
{\footnotesize\color{bkGray}接頭語：G（ギガ、$10^9$）、M（メガ、$10^6$）、k（キロ、$10^3$）、m（ミリ、$10^{-3}$）、μ（マイクロ、$10^{-6}$）、n（ナノ、$10^{-9}$）。\par}

%% ---------------- なぜ？の道すじ ----------------
\bkbacktitle{なぜ？の道すじ}
{\small 気になる「なぜ？」から、読む章の順番をたどる路線図です。札の横の数はこの巻のページ、ほかの巻の札（熱・波・電・原）は第2〜5巻の章です。使う数学の章は、下の小さな札です。（第13回の見本。50〜100の「なぜ？」を、物理と数学の両方のシリーズから集めて足す）\par}
\vspace{1mm}
\phwhyroute{エレベータが上へ動き出すと、体が重く感じるのはなぜ？}{力-2}{力-1,力-2}{4-11,5-5,6-2}
\phwhyroute{鉄でできた大きな船が、なぜ水に浮かぶ？}{力-12}{力-2,力-12}{5-12,8-6}
\phwhyroute{ジェットコースターは、最初の山より高く上がれる？}{力-4}{力-2,力-4}{4-12,8-7}
\phwhyroute{スケートの選手が腕をたたむと、速く回るのはなぜ？}{力-6}{力-3,力-6}{4-14,8-6}
\phwhyroute{本棚の板は、なぜまん中がいちばんたわむ？}{力-10}{力-5,力-8,力-10}{5-13,6-3}
\phwhyroute{ホースの先をつまむと、水が遠くまで飛ぶのはなぜ？}{力-13}{力-13,力-14}{8-8,6-3}
\phwhyroute{車はどうして動くの？}{力-0}{力-2,力-6,力-7,熱-7,電-12}{4-14,5-5,6-2}
\phwhyroute{冷たいコップのまわりに、なぜ水がつく？}{熱-9}{熱-1,熱-8,熱-9,熱-11}{3-5,3-11,6-4}
\phwhyroute{空はなぜ青い？　夕焼けはなぜ赤い？}{波-13}{波-5,波-13,電-15}{2-7,4-8,6-10}

%% ---------------- チェック表 ----------------
\bkbacktitle{できるようになったことチェック表}
{\small 各章の「ふりかえり」を集めた表です。章を学び終えたら1回め、次の巻に進む前に2回めのチェックをしよう。自信がない所は、ページの数字の所へもどろう。\par}\vspace{2mm}
% 自動生成（tools/mkchecklist.py p1）。各章の「ふりかえり」の項目を集めたもの（第13回）
\begingroup\footnotesize
\def\bkcb{\tikz[baseline=-0.5ex]\draw[\bkTF, line width=0.7pt, rounded corners=0.5pt] (0,-0.3em) rectangle (0.8em,0.5em);}
\begin{longtblr}{colspec={Q[c,m,wd=11mm] X[l,m] Q[c,m,wd=10mm] Q[c,m,wd=9mm] Q[c,m,wd=9mm]}, width=\linewidth, rowhead=1, colsep=3pt,
  row{1}={bg=bkVol!14, font=\sffamily\bfseries\scriptsize}, hlines={0.4pt,bkLine}, rowsep=0.7pt}
章 & できるようになったこと & ページ & 1回め & 2回め \\
― & （章ができたら、各章のふりかえりから自動で作る） & ― & \bkcb & \bkcb \\
\end{longtblr}
\endgroup


%% ---------------- つぎの巻へ ----------------
\bkbacktitle{つぎの巻へ}
\newcommand\bknextpill[1]{\bk@pill{#1}}
{\small 力学で、動く・回る・ゆがむ・流れる物を「数える量」で調べる道具がそろいました。第2巻『熱学』では、同じバランス式に「熱」を入れて、エンジン・冷蔵庫・発電所を調べます。\par}\vspace{3mm}
\begin{tcolorbox}[enhanced, colframe=phVol2, colback=white, boxrule=0.9pt, arc=2.5mm, left=3mm, right=3mm, top=2mm, bottom=2mm,
  title={第2巻　熱学：エネルギーを使いこなす}, fonttitle=\sffamily\bfseries, colbacktitle=phVol2, coltitle=white]
\small
\begin{itemize}[leftmargin=1.2\zw, itemsep=2pt, topsep=1pt, label={\color{phVol2}●}]
\item \bk@pill{熱-2}\ 仕事で温まる：第一法則（力-4 の仕事とエネルギーを、熱まで広げる）
\item \bk@pill{熱-4}\ 流れる気体の熱：エンタルピーとノズル（力-13 のベルヌーイの式とつながる）
\item \bk@pill{熱-12}\ 流れが熱を運ぶ：対流熱伝達（力-14・力-15 のレイノルズ数と境界層を使う）
\end{itemize}
\end{tcolorbox}
\phbkhead{この巻の道具は、この先でこう使う}
{\small この巻の章と、その道具を使う第2巻から先の章です（章のはじめの「道しるべ」の「この章の先に」から自動で集めた表）。\par}\vspace{2mm}
% 自動生成（tools/mknext.py p1）。各章の道しるべ「この章の先に」から、後ろの巻の章を集めたもの（第13回）
\begin{longtblr}{colspec={Q[l,m,wd=58mm] X[l,m]}, width=\linewidth, rowhead=1, row{1}={bg=bkVol!14, font=\sffamily\bfseries\footnotesize},
  hlines={0.4pt,bkLine}, rowsep=1.5pt, cells={font=\footnotesize}}
この巻の章 & この道具を使う章 \\
（章ができたら、道しるべから自動で作る） & ― \\
\end{longtblr}


%% ---------------- さくいん ----------------
\bkbacktitle{さくいん}
{\small 数字はページです。その語がはじめて出てくるページや、くわしく説明しているページをのせています。}\par\vspace{2mm}
\begin{multicols}{2}\small\setlength{\parskip}{0pt}
\directlua{bk.printidx()}
\end{multicols}

%% ---------------- 奥付 ----------------
\bkpadfour
\global\bk@decofalse\bkfolio{false}
\null\vfill
\noindent\begin{minipage}{\linewidth}\small
{\sffamily\bfseries\Large\color{bkDeep}\phSeries　第1巻}\par\vspace{1mm}
{\sffamily\bfseries\large 力学：動く・回る・ゆがむ・流れる}\par\vspace{3mm}
{\color{bkVol}\hrule height 1pt}\vspace{3mm}
\begin{tabular}{@{}l@{\hspace{6mm}}l@{}}
編者 & \phEditor \\
 & 委員長　小学校長（まとめ役） \\
 & 委員　数学・電気工学・機械工学・教育学・量子力学・建築工学・\\
 & \phantom{委員　}半導体工学の大学教授、幼稚園教諭、塾講師、\\
 & \phantom{委員　}特別支援学校教諭、心理士、科学解説者、小学校教諭 \\
デザイン & ユニバーサルデザイン担当委員 \\
挿絵 & イラスト担当委員 \\
先に読む本 & 「物理を理解するための数学」全9巻 \\
版 & 第1版 Ver1（第13回編集委員会）　2026年10月 \\
\end{tabular}\par\vspace{3mm}
{\color{bkVol}\hrule height 1pt}\vspace{2mm}
{\footnotesize この本の問題の答えは、巻末の「こたえ」にあります。くわしい解き方は、別冊の問題集と指導書にのせる予定です。}
\end{minipage}
\makeatother
